\section{Paper Collection Protocol}
\label{app:method}

The search is scripted and rerunnable; the configuration (including every query string and term list), raw hits, per-stage counts, and the full-text coding record for the runs reported here (identifiers \texttt{20260915\_v1} for 2025--2026 and \texttt{20260916\_v3} for the 2020--2026 window) are released with the paper.

\paragraph{Queries.}
Twenty-three keyword queries were defined in the first run, four to five per survey section; a second run added eleven queries on evidence acquisition and theory-of-mind belief modelling (34 in total): representation (Section~\ref{sec:representation}), revision (Section~\ref{sec:revision}), learning (Section~\ref{sec:learning}), evidence (Section~\ref{sec:evidence}), and evaluation (Section~\ref{sec:evaluation}).

\paragraph{Sources.}
Source A queried OpenAlex across all venues with publication date from 1~January 2025, 50 results per query. Source B queried OpenAlex restricted to the arXiv source with publication date from 1~January 2026, 50 results per query, to improve recall of preprints. The arXiv API itself was rate-limited during the run; the OpenAlex mirror of arXiv lags by a few days.

\paragraph{Screening.}
Records were deduplicated by identifier and then by normalised title, so that a preprint and its journal version count once. Rule-based screening on title and abstract applied one inclusion criterion for an LLM-related term (I1), one exclusion criterion for off-domain applications (E1), and one inclusion criterion requiring at least two topic terms from a fixed list (I2). Records without an abstract of at least 200 characters were excluded. Included records were ranked by a relevance score used only for ordering the manual full-text screen, never for inclusion.

\begin{table}[h]
\centering\small
\begin{tabular}{lr}
\toprule
Stage & Count \\
\midrule
Queries & 23 \\
Source A raw hits / unique & 1{,}150 / 715 \\
Source B raw hits / unique & 1{,}150 / 1{,}008 \\
Merged unique (identifier) & 1{,}685 \\
After title-level deduplication & 1{,}460 \\
Excluded: no or short abstract & 45 \\
Excluded I1: no LLM term & 160 \\
Excluded E1: off-domain & 35 \\
Excluded I2: fewer than two topic terms & 357 \\
Included after rule screening & 863 \\
\quad of which published in 2026 & 656 \\
Full-text screened & 117 \\
\quad of which included after reading & 106 \\
\bottomrule
\end{tabular}
\caption{Paper-collection counts for run \texttt{20260915\_v1}.}
\label{tab:counts}
\end{table}

\paragraph{Second and third runs.}
The second run (34 queries, same sources and rules) returned 834 unique records, 338 after screening, of which 96 were published in 2026; only two of those had not already been read in the first run, which we take as evidence that the first run had saturated the keyword neighbourhood. The third run extended the window to 1~January 2020 with 100 results per query and returned 4{,}471 unique records, 3{,}999 after title deduplication and 2{,}079 after rule screening. Keyword ranking is dominated by generic high-citation papers before 2025, so the 2020--2025 foundations were enumerated as a curated list of 90 titles resolved through OpenAlex for venue and citation data (Table~\ref{tab:masterfull}); 75 of them were read in full and 73 included.

\paragraph{Citation-tree expansion.}
Keyword search cannot reach papers that describe belief maintenance in other words, so the included set was expanded by snowballing over the citation graph. Each round took the papers newly included in the previous round as seeds, fetched their reference lists (backward) and the papers citing them (forward) from the Semantic Scholar Graph API, kept candidates published from 2020, removed everything already screened, and applied the same rule-based screen as the keyword runs. Because forward citation of recent preprints is dense, the rule-included candidates were sampled for full-text reading by a fixed rule (topic score $\geq 10$, or the word \emph{belief} in title or abstract, or linked to at least three seeds, or a world-model paper from 2025 with score $\geq 7$), and every sampled paper was read and coded with the same fields as the keyword runs. Round one was seeded by the 179 keyword-included papers and each later round by the papers the previous round added; the full-text sample shrank from 302 to 115, 83, and 26 across four rounds, and the expansion stopped after the round whose sample fell below thirty, the stopping rule fixed in advance (Table~\ref{tab:snowball}).

\begin{table}[h]
\centering\footnotesize
\setlength{\tabcolsep}{3pt}
\resizebox{\columnwidth}{!}{%
\begin{tabular}{lrrrr}
\toprule
Stage & R1 & R2 & R3 & R4 \\
\midrule
Seeds & 178 & 165 & 55 & 45 \\
Backward references & 9{,}127 & 8{,}601 & 2{,}673 & 2{,}406 \\
Forward citations & 9{,}999 & 9{,}999 & 1{,}403 & 441 \\
Unique candidates (2020--) & 12{,}872 & 13{,}113 & 2{,}536 & 1{,}465 \\
Already screened & 302 & 1{,}284 & 605 & 465 \\
Excluded by rules & 8{,}448 & 7{,}401 & 1{,}300 & 747 \\
Included after rule screening & 4{,}122 & 4{,}428 & 631 & 253 \\
Sampled for full-text reading & 302 & 115 & 83 & 26 \\
Included after reading & 171 & 58 & 47 & 7 \\
\quad of which peer-reviewed & 79 & 23 & 22 & 2 \\
\bottomrule
\end{tabular}}
\caption{Citation-tree expansion by round. Forward citations are capped at 9{,}999 per request batch by the API.}
\label{tab:snowball}
\end{table}

\paragraph{Year distribution.}
Table~\ref{tab:years} gives the included papers by year, source, and evidence level. The share that is peer-reviewed falls with recency, from four in five before 2025 to about one in nine for 2026, which is the ordinary lag of review rather than a property of the field; the evidence labels in every table make the dependence explicit.

\begin{table}[h]
\centering\footnotesize
\setlength{\tabcolsep}{4pt}
\begin{tabular}{lrrrrrrr}
\toprule
Year & Keyword & R1 & R2 & R3 & R4 & Total & Peer-rev. \\
\midrule
2020 & 1 & 3 & 1 & 0 & 0 & 5 & 4 \\
2021 & 2 & 2 & 1 & 0 & 0 & 5 & 3 \\
2022 & 9 & 4 & 0 & 1 & 0 & 14 & 13 \\
2023 & 27 & 21 & 2 & 0 & 0 & 50 & 41 \\
2024 & 16 & 20 & 9 & 4 & 1 & 50 & 40 \\
2025 & 18 & 69 & 20 & 20 & 2 & 129 & 49 \\
2026 & 106 & 52 & 25 & 22 & 4 & 209 & 23 \\
\midrule
Total & 179 & 171 & 58 & 47 & 7 & 462 & 173 \\
\bottomrule
\end{tabular}
\caption{Included papers by publication year and by the run that found them (keyword search or citation-tree round).}
\label{tab:years}
\end{table}

\paragraph{Evidence level.}
Each cited source is labelled P (peer-reviewed venue) or A (arXiv preprint at the time of writing) in the tables and in Figure~\ref{fig:metrics}. Of the 462 papers included after full-text reading (179 from keyword search and 283 from citation-tree expansion), 173 are peer-reviewed and 289 are preprints; every number quoted in the text is the paper's own headline figure, taken from the abstract or the reported table, and none has been reproduced by us.

\paragraph{Full-text screening.}
A paper passed full-text screening if it (a) concerns a language-model agent or a component used by one, and (b) constructs, revises, tests, learns, or evaluates an explicit state carried between decisions, or documents a failure of such a state; benchmarks and surveys that do (b) were kept and coded as such. Each paper read in full was recorded with the fields: who maintains the belief (source of the state term), credited object, revision signal (source of the transition and likelihood terms), whether belief-level metrics are reported, survey section, and role. The record is the source of Figure~\ref{fig:taxonomy}, Tables~\ref{tab:representation}--\ref{tab:benchmarks}, and Figure~\ref{fig:metrics}. Coding was done by the authors; an independent re-coding of a subset is planned and not yet available.

\paragraph{Peer-reviewed subset.}
Because 289 of the 462 papers are preprints, Table~\ref{tab:sensitivity} recomputes the counts behind each insight box on the 173 peer-reviewed papers alone.

\begin{table}[h]
\centering\scriptsize
\setlength{\tabcolsep}{3pt}
\begin{tabular}{@{}clrr@{}}
\toprule
\# & Counts the insight rests on & All & P only \\
\midrule
1 & state term written / filter / world model & 66 / 20 / 54 & 22 / 7 / 19 \\
2 & revision papers (\S\ref{sec:revision}) & 106 & 34 \\
3 & credited: interaction / state / state by own obs. & 93 / 47 / 1 & 32 / 12 / 0 \\
4 & acquisition papers (\S\ref{sec:evidence}) & 61 & 18 \\
5 & papers with a belief-level metric & 88 & 33 \\
\bottomrule
\end{tabular}
\caption{Sensitivity of the five insights to preprint reliance: the counts each rests on, over all included papers and over the peer-reviewed subset alone. Every insight is a claim about a proportion or an absence, and each holds in the same direction on the peer-reviewed subset; the cell for state credit signed by the agent's own observation holds one paper, a preprint, in both.}
\label{tab:sensitivity}
\end{table}

